% Source: https://artofproblemsolving.com/community/c3h2669582_issue_with_djmathman_100_geometry_problems
% Stable identity: aops-topic-2669582-q01
Use complex numbers with $P=0$.

\textbf{(a)} A similarity preserving orientation and fixing $P=0$ is $z\mapsto\lambda z$ for a nonzero complex number $\lambda$ (a rotation about $P$ followed by a scaling). So $C=\lambda A$ and $D=\lambda B$. Let $\mu=\dfrac BA$, so $B=\mu A$. Then
\[D=\lambda B=\lambda\mu A=\mu(\lambda A)=\mu C.\]
So the map $z\mapsto\mu z$, also a rotation and scaling about $P$, sends $P,A,C$ to $P,B,D$. Hence $\triangle PAC\sim\triangle PBD$. $\blacksquare$

\begin{asy}
usepackage("amsmath");
size(4cm);
defaultpen(fontsize(10pt));
pair P=(0,0),A=(2,0.3),B=(1.4,1.5);real k=1.5,th=50;
pair C=k*(rotate(th)*A),D=k*(rotate(th)*B);
draw(P--A--B--cycle,linewidth(0.8));draw(P--C--D--cycle,linewidth(0.8));
draw(A--C,dashed+blue);draw(B--D,dashed+blue);
dot(P);dot(A);dot(B);dot(C);dot(D);
label("$P$",P,SW);label("$A$",A,SE);label("$B$",B,N);label("$C$",C,E);label("$D$",D,N);
\end{asy}

\textbf{(b)} Take $P=0$, $A=1$, $B=1+2i$, and reflect in the line $y=x$, i.e. apply $z\mapsto i\,\overline z$: $C=i$ and $D=2+i$. Triangle $PCD$ is the mirror image of triangle $PAB$, so the two are similar (even congruent) with opposite orientations. But triangle $PAC$ has sides $PA=PC=1$, $AC=\sqrt2$, while triangle $PBD$ has sides $PB=PD=\sqrt5$, $BD=\sqrt2$. Both are isosceles with apex $P$, and $\dfrac{AC}{PA}=\sqrt2\ne\dfrac{BD}{PB}=\dfrac{\sqrt2}{\sqrt5}$, so they are not similar. $\blacksquare$

\begin{asy}
usepackage("amsmath");
size(3.5cm);
defaultpen(fontsize(10pt));
pair P=(0,0),A=(1,0),B=(1,2),C=(0,1),D=(2,1);
draw((-0.3,-0.3)--(2.3,2.3),gray(0.7)+dotted);
draw(P--A--B--cycle,linewidth(0.8));draw(P--C--D--cycle,linewidth(0.8));
draw(A--C,dashed+blue);draw(B--D,dashed+blue);
dot(P);dot(A);dot(B);dot(C);dot(D);
label("$P$",P,SW);label("$A$",A,SE);label("$B$",B,N);label("$C$",C,W);label("$D$",D,E);label("$y=x$",(2.3,2.3),NE,fontsize(8pt));
\end{asy}
